\clearpage\begingroup\SourceLeaf{35}\HistoricalHead{30}{L}{PHYSIK. BRIEFWECHSEL.}
\begin{center}
[7.]\\[.7em]
\textsc{Gauss} an \textsc{Schumacher}. Göttingen, 24. Julius 1836.
\par\medskip\rule{.78\textwidth}{.35pt}\par
{\small [Briefwechsel zwischen \textsc{Gauss} und \textsc{Schumacher} III, Altona 1860, S. 100.]}\par
\rule{.78\textwidth}{.35pt}
\end{center}
\noindent .\ .\ .\ .\ .\ .\ .\ .\ .\ .\ .\ .\ .\ .\ .\ .\ .\ .\ .\ .\ .
\par\noindent\hspace*{\parindent}Eine Hauptabweichung bei meiner Wägungsart von der gewöhnlichen
ist, dass ich \textsc{Bordas} Manier habe fahren lassen. Ich begreife in der That
nicht, warum man sich an diese gehalten hat, in den Fällen, wo man die
grösste Genauigkeit verlangt, also \textls[120]{oft} wiederholte Wägungen macht. Es ist
sehr klar, dass wenn Sie gar nicht tariren, sondern die beiden zu vergleichenden
Gewichte auf den Schalen umtauschen, ihre Ungleichheit einen doppelt
so grossen Ausschlag gibt, wie bei \textsc{Bordas} Art, und dass man also mit einer
gegebenen Waage bei meiner Wägungsart mit 20 Wägungen gerade eben so
weit kommt, wie mit 80 Wägungen nach \textsc{Bordas} Art (eine doppelt so scharfe
Operation hat nemlich bekanntlich das 4fache Gewicht). Welch enormer Gewinn!
Freilich müssen noch einige Nebenumstände dabei berücksichtigt werden
(die zu erwähnen hier zu. weitläuftig sein würde), die aber nicht die geringste
Schwierigkeit haben. Bei 32 Wägungen des {\fontspec{FreeSerif.otf}℔}, die ich nach \textls[120]{dieser} Art
gemacht habe, war der mittlere Fehler \textls[120]{Einer} Wägung nur 0,2 Milligramm,
und es ist mir noch keine einzige vorgekommen, die ein halbes Milligramm
vom Mittel differirt hätte.\quad .\ .\ .\ .\ .\ .\ .
\clearpage\endgroup
